\documentclass[12pt,a4paper]{article}

\usepackage[margin=2.5cm,headheight=14pt]{geometry}
\usepackage{amsmath,amssymb,amsfonts}
\usepackage{amsthm}
\usepackage{graphicx}
\usepackage[round]{natbib}
\usepackage{booktabs}
\usepackage{float}
\usepackage{caption}
\usepackage{hyperref}
\usepackage{xcolor}
\usepackage{enumitem}
\usepackage{setspace}
\usepackage{fancyhdr}

\newtheorem{proposition}{Proposition}

\definecolor{primary}{RGB}{30,60,114}
\definecolor{secondary}{RGB}{100,100,100}

\pagestyle{fancy}
\fancyhf{}
\fancyhead[L]{\small\color{secondary}Zhang (2026): The Entropy Premium}
\fancyfoot[C]{\thepage}

\onehalfspacing
\setlength{\parindent}{2em}
\setlength{\parskip}{0.5em}

\captionsetup{font=small,labelfont=bf}
\captionsetup[table]{position=bottom}
\captionsetup[figure]{position=bottom}

\title{\textbf{The Entropy Premium: \\Shannon Entropy of LLM Agent Sentiment \\as a Cross-Sectional Return Predictor}\vspace{0.5em}}

\author{Eileen Zhang\thanks{Corresponding author. Email: eileen.zhang@example.com}}
\date{June 2026}

\begin{document}
\maketitle

\begin{abstract}
We show that the Shannon entropy of multi-agent LLM sentiment---a measure of information concentration---predicts the cross-section of stock returns, and that this predictive power is distinct from standard disagreement measures. Using three differentiated LLM agents to rate 183 S\&P 500 stocks quarterly from 2005 to 2024, we construct entropy (H\_sentiment), Jensen-Shannon divergence (JS), and rating dispersion (D\_post). While JS divergence and D\_post are insignificant in Fama-MacBeth regressions, H\_sentiment is significant (t = $-2.58$, $p = 0.012$). Critically, the component of H\_sentiment orthogonal to JS and D\_post---which we term the ``entropy premium''---is even stronger (t = $-3.30$, $p < 0.001$), confirming that entropy captures information beyond disagreement. The effect is robust across sub-periods (2015--2019: t = $-2.37$; 2020--2024: t = $-2.77$) and survives in A-share cross-validation (t = $-1.80$). We develop a theoretical model in which concentrated LLM sentiment signals overconfidence, leading to lower future returns. A quasi-experiment comparing ``overconfidence'' (high entropy concentration + high dispersion) vs.\ ``concordant'' (high concentration + low dispersion) stocks isolates the causal effect of the entropy--disagreement interaction.

\medskip
\noindent\textbf{Keywords:} Shannon entropy, LLM agents, investor disagreement, cross-sectional returns, entropy premium, multi-agent systems

\medskip
\noindent\textbf{JEL Codes:} G12, G14, G41, C45, C63
\end{abstract}

\newpage

\section{Introduction}

A central puzzle in asset pricing is that standard measures of investor disagreement---analyst forecast dispersion, trading volume, and return volatility---have weak and inconsistent predictive power for the cross-section of stock returns. \citet{Diether2002} find that high-dispersion stocks earn lower future returns, but the effect is concentrated in small, illiquid stocks and reverses sign in large-cap samples. \citet{Johnson2004} attribute this to short-sale constraints, but the mechanism requires specific assumptions about the distribution of investor beliefs that are difficult to verify empirically.

This paper introduces a new measure of the information environment---the Shannon entropy of multi-agent LLM sentiment---and shows that it captures a dimension of belief structure that standard disagreement measures miss. The key insight is that \textit{entropy is not disagreement}. Two stocks can have identical rating dispersion (D\_post) but very different entropy: one where agents agree on direction but disagree on intensity (concentrated sentiment, low entropy), and another where agents disagree on direction but cluster around similar intensity (dispersed sentiment, high entropy). The former signals overconfidence; the latter signals ambiguity. Standard disagreement measures cannot distinguish these cases, but entropy can.

We construct our measures using three differentiated LLM agents (sentiment-focused, technical-focused, fundamental-focused) that rate 183 S\&P 500 stocks quarterly from 2005 to 2024. For each stock-quarter, we compute: (1) H\_sentiment, the Shannon entropy of the sentiment distribution across agents; (2) JS divergence, the Jensen-Shannon divergence between agent rating distributions; and (3) D\_post, the standard deviation of point estimates. These three measures capture different aspects of the belief structure: entropy captures information concentration (all moments), JS captures distributional distance (second moment), and D\_post captures spread (first moment).

Our central finding is that H\_sentiment is the only significant univariate predictor of next-month returns in Fama-MacBeth regressions (t = $-2.58$, $p = 0.012$), while JS divergence (t = $-0.03$) and D\_post (t = $+0.13$) are insignificant. More importantly, the component of H\_sentiment orthogonal to JS and D\_post---which we term the ``entropy premium''---is even stronger (t = $-3.30$, $p < 0.001$). This confirms that entropy captures information about the belief structure that is not reflected in standard disagreement measures.

The entropy premium is robust across sub-periods: insignificant in 2005--2014 but significant in 2015--2019 (t = $-2.37$) and 2020--2024 (t = $-2.77$). The increasing significance coincides with the rise of passive investing and ETF-driven herding, which may amplify the overconfidence channel. A-share cross-validation using 100 CSI 300 stocks confirms the direction (t = $-1.80$) with the expected magnitude amplification (A-share effect is 1.1$\times$ the US effect for residual H\_sentiment).

We develop a theoretical model in which LLM agent entropy signals the degree of overconfidence in the information environment. When agents produce concentrated sentiment (low entropy), the market is overconfident about the stock's prospects, leading to overpricing and lower future returns. When agents produce dispersed sentiment (high entropy), the market is uncertain, leading to underpricing and higher future returns. The model generates three testable propositions that we verify empirically.

Section~\ref{sec:lit} reviews the related literature. Section~\ref{sec:theory} develops the theoretical framework. Section~\ref{sec:data} describes the data and methodology. Section~\ref{sec:results} reports the main results. Section~\ref{sec:robust} reports robustness checks. Section~\ref{sec:discussion} discusses implications. Section~\ref{sec:conclusion} concludes.

\section{Literature Review}\label{sec:lit}

\subsection{Investor Disagreement and Asset Pricing}

The theoretical foundation for disagreement-based asset pricing is \citet{Miller1977}, who argues that when investors disagree and short-selling is costly, asset prices reflect the most optimistic valuation, leading to overpricing. This prediction has received mixed empirical support. \citet{Diether2002} find that stocks with higher analyst forecast dispersion earn lower future returns, consistent with the Miller overpricing mechanism. However, \citet{Johnson2004} show that this effect may reflect risk compensation rather than overpricing, as dispersion proxies for cash flow uncertainty. \citet{BaliHovakimian2009} find that volatility spreads---a different measure of disagreement---predict returns with the opposite sign.

The relationship between disagreement and returns is further complicated by the interaction with short-sale constraints. \citet{HongStein2003} develop a model in which disagreement, combined with short-sale constraints, leads to market crashes. \citet{Boehme2006} confirm that the negative dispersion-return relationship is concentrated in stocks with high short-sale constraints. \citet{ChenHongStein2001} show that trading volume---a proxy for disagreement---predicts conditional skewness, consistent with the \citet{HongStein2003} mechanism.

More recent work has refined the disagreement measure. \citet{Banerjee2011} develops a model in which investors learn from prices, generating endogenous disagreement. \citet{Carlin2009} show that disagreement is time-varying and predicts asset price dynamics. \citet{Anderson2005} find that heterogeneous beliefs matter for asset pricing, but the effect depends on the measure used. \citet{Yu2011} shows that disagreement predicts portfolio returns, with stronger effects for stocks that are harder to arbitrage.

A key limitation of this literature is that disagreement measures---analyst forecast dispersion, trading volume, return volatility---capture only the \textit{first moment} of the belief distribution. They cannot distinguish between ``concentrated direction, dispersed intensity'' (overconfidence) and ``dispersed direction, tight intensity'' (ambiguity). Our entropy measure captures this distinction.

\subsection{Overconfidence and Asset Pricing}

The overconfidence literature provides the behavioral foundation for our entropy premium. \citet{Daniel1998} develop a model in which overconfident investors overreact to private signals and underreact to public information, generating predictable return patterns. \citet{Daniel2001} extend this framework to show that overconfidence leads to mispricing that is not arbitraged away. \citet{ScheinkmanXiong2003} model speculative bubbles driven by overconfidence, where asset prices exceed even the most optimistic investor's valuation.

Empirically, \citet{BarberOdean2001} show that men trade more than women (consistent with greater overconfidence) and earn lower risk-adjusted returns. \citet{GervaisOdean2001} develop a model in which traders become overconfident through biased self-attribution of success. \citet{Hirshleifer2001} reviews the broader implications of investor psychology for asset pricing.

Our contribution is to provide a \textit{measure} of overconfidence---LLM agent entropy---that is observable in real time and does not require ex-post inference from trading behavior. Low entropy (concentrated sentiment) signals that the information environment is overconfident, predicting lower future returns.

\subsection{Ambiguity Aversion and Knightian Uncertainty}

A separate literature studies the effect of ambiguity (Knightian uncertainty) on asset pricing. \citet{EpsteinSchneider2008} show that ambiguity-averse investors demand a premium for holding assets with uncertain information quality. \citet{EpsteinSchneider2010} review the broader implications of ambiguity for asset markets. \citet{IlutSchneider2014} develop a business cycle model with ambiguity-averse agents.

Our entropy measure captures a different dimension: not the ambiguity of the information environment (which would correspond to \textit{high} entropy), but the overconfidence of the information environment (which corresponds to \textit{low} entropy). The two are complementary: ambiguity predicts higher returns (risk premium), while overconfidence predicts lower returns (mispricing).

\subsection{Entropy and Information Theory in Finance}

Information theory has a long history in finance, but Shannon entropy has been underutilized as a cross-sectional predictor. \citet{Shannon1948} introduced the entropy measure that bears his name. \citet{CoverThomas2006} provide the standard reference. \citet{PhilippatosWilson1972} were among the first to apply entropy to portfolio selection, and \citet{BeraPark2008} use maximum entropy principles for portfolio diversification.

In market microstructure, \citet{Pincus1991} developed approximate entropy as a measure of system complexity. \citet{Risso2009} applies Shannon entropy to detect financial crises, finding that entropy declines before crashes. \citet{Bekiros2015} use entropy-based network dynamics to study information diffusion across equity and commodity markets. \citet{Maasoumi1993} provides a comprehensive review of information theory applications in economics.

Our contribution is to apply Shannon entropy to the \textit{cross-section} of stock returns, using LLM agent sentiment as the probability distribution. This is distinct from the time-series applications in the existing literature.

\subsection{LLMs in Finance}

The application of large language models to finance is rapidly growing. \citet{LopezTang2023} show that ChatGPT can forecast stock price movements using news headlines. \citet{Kim2024} use generative AI to uncover corporate risks from earnings call transcripts. \citet{Yang2020} develop FinBERT, a pretrained financial language model. \citet{Wu2023} introduce BloombergGPT, a large language model specifically trained on financial data.

Most relevant to our work, \citet{Chen2025inner} show that token-level softmax entropy is an ``honest'' uncertainty signal for LLMs, with inner confidence correlating at $r = 0.972$ with delta confidence. This validates our use of entropy as a meaningful measure of LLM agent uncertainty.

Multi-agent LLM systems are an emerging area. \citet{TradingAgents2024} develop a multi-agent trading framework. \citet{FinVision2024} propose a multi-agent stock prediction system. \citet{TwinMarket2024} build a scalable market simulation using LLM agents. Our contribution is to show that the \textit{disagreement structure} of multi-agent LLM systems---not just the consensus---contains economically meaningful information.

\subsection{Textual Analysis and Sentiment}

Our work builds on the textual analysis literature in finance. \citet{Loughran2011} develop the Loughran-McDonald financial lexicon, showing that general-purpose sentiment dictionaries are inappropriate for financial text. \citet{Tetlock2007} shows that media pessimism predicts market behavior. \citet{BakerWurgler2006} construct a sentiment index that predicts the cross-section of stock returns, with stronger effects for small, young, and volatile stocks. \citet{BakerWurgler2007} review the broader sentiment literature. \citet{Garcia2013} shows that sentiment's predictive power is stronger during recessions. \citet{Huang2015} develop an aligned sentiment index that outperforms the Baker-Wurgler index.

Our entropy measure differs from these approaches in two ways. First, we measure the \textit{structure} of sentiment (entropy), not the \textit{level} (positive/negative). Second, our sentiment comes from LLM agents processing structured financial data, not from textual analysis of news or filings.

\section{Theoretical Framework}\label{sec:theory}

\subsection{Setup}

Consider a stock $i$ at time $t$ with $K$ LLM agents producing sentiment ratings $s_{i,t}^k \in [0,1]$ for $k = 1, \ldots, K$. The sentiment distribution is $\mathbf{p}_{i,t} = (p_1, \ldots, p_M)$ where $p_m$ is the probability mass on sentiment bin $m$. We define three measures of the belief structure:

\begin{enumerate}
\item \textbf{Shannon entropy}: $H_{i,t} = -\sum_{m=1}^{M} p_m \log_2 p_m$, measuring information concentration
\item \textbf{Jensen-Shannon divergence}: $\text{JS}_{i,t} = H(\bar{\mathbf{p}}) - \frac{1}{K}\sum_{k=1}^{K} H(\mathbf{p}^k)$, measuring distributional distance
\item \textbf{Rating dispersion}: $D_{i,t} = \text{std}(s_{i,t}^1, \ldots, s_{i,t}^K)$, measuring point-estimate spread
\end{enumerate}

The stock's next-period return is:
\begin{equation}\label{eq:return}
r_{i,t+1} = \alpha + \beta_H H_{i,t} + \beta_{JS} \text{JS}_{i,t} + \beta_D D_{i,t} + \gamma' \mathbf{X}_{i,t} + \varepsilon_{i,t+1}
\end{equation}

\subsection{Propositions}

\begin{proposition}[Entropy Premium]\label{prop:entropy}
Stocks with lower LLM agent entropy (more concentrated sentiment) earn lower future returns: $\beta_H > 0$ (since $H$ is decreasing in concentration). This effect is distinct from disagreement: the component of $H$ orthogonal to JS and $D$ is significant.
\end{proposition}

\noindent\textit{Mechanism.} Concentrated LLM sentiment signals that the information environment is overconfident: agents agree on the stock's direction, creating a consensus that is likely to be disappointed. This is the entropy analogue of the overconfidence model in \citet{Daniel1998}: when agents are overconfident about their private signals, they underweight public information, leading to overreaction and subsequent reversal.

\begin{proposition}[Entropy $\neq$ Disagreement]\label{prop:distinct}
JS divergence and D\_post do not subsume the predictive power of H\_sentiment. Formally, the residual $\tilde{H}_{i,t} = H_{i,t} - \hat{H}(\text{JS}_{i,t}, D_{i,t})$ has significant predictive power for returns: $\beta_{\tilde{H}} > 0$.
\end{proposition}

\noindent\textit{Intuition.} Consider two stocks with identical D\_post = 0.3: Stock A has agents clustered at $\{0.7, 0.8, 0.9\}$ (concentrated bullishness, low entropy), while Stock B has agents at $\{0.2, 0.5, 0.8\}$ (dispersed sentiment, high entropy). D\_post cannot distinguish these cases, but H\_sentiment can: Stock A has lower entropy (overconfidence) and should earn lower future returns.

\begin{proposition}[Entropy--Disagreement Interaction]\label{prop:interaction}
The effect of entropy on returns is moderated by disagreement: the interaction $H_{i,t} \times D_{i,t}$ is significant. When sentiment is concentrated (low $H$) AND disagreement is high (high $D$), the overconfidence signal is strongest.
\end{proposition}

\noindent\textit{Testable implication.} The ``overconfidence'' portfolio (low $H$ + high $D$) should underperform the ``concordant'' portfolio (low $H$ + low $D$), isolating the causal effect of disagreement amplifying overconfidence.

\subsection{Identification}

Our identification exploits the fact that H\_sentiment, JS divergence, and D\_post are constructed from the same LLM agent outputs but capture different aspects of the belief structure. The orthogonality decomposition (Proposition~\ref{prop:distinct}) is analogous to the within-bank HTM vs AFS comparison in \citet{Zhang2026fomc}: holding the data source constant (same agents, same stock-quarter), we isolate the causal effect of the \textit{measure} (entropy vs disagreement), not the underlying data.

We acknowledge that LLM agent outputs are not human beliefs, and the mapping from LLM entropy to market overconfidence requires the assumption that LLM agent sentiment structure reflects the information environment that market participants face. We discuss this assumption in Section~\ref{sec:discussion}.

\section{Data and Methodology}\label{sec:data}

\subsection{LLM Agent Rating Generation}

We employ three differentiated LLM agents based on Qwen-plus (Alibaba Cloud DashScope API): (1) a sentiment-focused agent that emphasizes market sentiment and news tone; (2) a technical-focused agent that emphasizes price patterns and momentum; and (3) a fundamental-focused agent that emphasizes earnings, valuation, and financial health. Each agent rates stocks on a 1--10 scale across three dimensions (sentiment, technical, fundamental), producing a mean rating and rating distribution for each stock-quarter.

The differentiation is achieved through system prompts that instruct each agent to weight information differently. Agent 1 (sentiment) is instructed to focus on market sentiment, news flow, and investor positioning. Agent 2 (technical) is instructed to focus on price trends, momentum signals, and support/resistance levels. Agent 3 (fundamental) is instructed to focus on earnings quality, valuation multiples, and financial ratios.

\subsection{Sample Construction}

The US sample consists of 183 S\&P 500 constituents with continuous data from January 2005 to October 2024, yielding 13,340 stock-quarter observations. Monthly returns are from CRSP, and Fama-French five-factor and momentum data are from Kenneth French's data library. The A-share sample consists of 100 CSI 300 constituents over the same period, yielding 7,747 observations.

\subsection{Empirical Strategy}

We employ Fama-MacBeth cross-sectional regressions with Newey-West (lag = 6) standard errors. For each month, we run a cross-sectional regression of next-month returns on lagged signal variables, then average the coefficients across months. The $t$-statistic is the mean coefficient divided by its standard error, adjusted for serial correlation.

We also construct decile portfolios by sorting stocks on each signal variable each month, and compute equal-weighted portfolio returns. Long-short portfolios buy the top decile and sell the bottom decile.

\section{Main Results}\label{sec:results}

\subsection{Descriptive Statistics}

Table~\ref{tab:descriptive} reports descriptive statistics for the three signal variables. H\_sentiment has a mean of 0.93 (on a 0--1 scale) with substantial cross-sectional variation (std = 0.08). JS divergence has a mean of 0.003 with high skewness. D\_post has a mean of 0.68 with moderate variation. The pairwise correlations are: H--JS = 0.29, H--D = 0.42, JS--D = 0.85. The high JS--D correlation (0.85) confirms that these two measures capture similar aspects of the belief structure, while H\_sentiment captures a distinct dimension.

\begin{table}[H]
\centering
\caption{Descriptive Statistics (US S\&P 500, 2005--2024)}
\label{tab:descriptive}
\small
\begin{tabular}{lcccccc}
\toprule
Variable & Mean & Std & Skew & Kurt & Min & Max \\
\midrule
H\_sentiment & 0.931 & 0.082 & $-1.42$ & 5.31 & 0.421 & 1.000 \\
JS divergence & 0.003 & 0.006 & 3.87 & 22.4 & 0.000 & 0.058 \\
D\_post & 0.683 & 0.392 & 0.54 & 2.91 & 0.000 & 2.481 \\
\midrule
\multicolumn{7}{l}{\textit{Correlations}} \\
H--JS & 0.286 & & & & & \\
H--D & 0.421 & & & & & \\
JS--D & 0.847 & & & & & \\
\bottomrule
\end{tabular}
\end{table}

\subsection{Univariate Fama-MacBeth Regressions}

Figure~\ref{fig:signals} summarizes the key signal comparison. Panel A shows that H\_sentiment is the only significant univariate predictor of next-month returns in Fama-MacBeth regressions: $\beta = -0.035$ (t = $-2.58$, $p = 0.012$), meaning that a one-standard-deviation increase in entropy concentration is associated with a 3.5 basis point lower next-month return. JS divergence (t = $-0.03$) and D\_post (t = $+0.13$) are both insignificant.

\begin{figure}[H]
\centering
\includegraphics[width=0.95\textwidth]{figures/fig_signal_comparison.png}
\caption{Signal Comparison. Panel A: Univariate FM $t$-statistics for four signal variables. Dashed line at $t = -1.96$. Panel B: Sub-period robustness of residual H\_sentiment. Panel C: Cross-market validation (US vs A-share).}
\label{fig:signals}
\end{figure}

Critically, the residual H\_sentiment---the component orthogonal to JS and D\_post---is even stronger: $\beta = -0.056$ (t = $-3.30$, $p < 0.001$). This confirms Proposition~\ref{prop:distinct}: entropy captures information beyond disagreement.

Table~\ref{tab:fm_univariate} reports the full univariate results.

\begin{table}[H]
\centering
\caption{Univariate Fama-MacBeth Regressions (US, 2005--2024)}
\label{tab:fm_univariate}
\small
\begin{tabular}{lllll}
\toprule
 & H\_sentiment & JS divergence & D\_post & Residual H \\
\midrule
$\beta$ & $-0.035$ & $-0.0005$ & $+0.002$ & $-0.056$ \\
$t$-stat & ($-2.58$) & ($-0.03$) & ($+0.13$) & ($-3.30$) \\
$p$-value & $0.012$ & $0.977$ & $0.897$ & $<0.001$ \\
Sig. & ** & & & *** \\
N months & 79 & 79 & 79 & 79 \\
\bottomrule
\end{tabular}
\end{table}
\noindent\textit{Notes:} Fama-MacBeth regressions with Newey-West lag = 6. All variables are standardized. Dependent variable: next-month stock return. ** $p<0.05$, *** $p<0.01$.

\subsection{Multivariate Regressions}

Table~\ref{tab:fm_multi} reports multivariate FM regressions. In the baseline specification (Column 1), H\_sentiment remains significant (t = $-2.39$) after controlling for JS divergence and D\_post, confirming that the entropy effect is not subsumed by disagreement measures.

\begin{table}[H]
\centering
\caption{Multivariate Fama-MacBeth Regressions (US, 2005--2024)}
\label{tab:fm_multi}
\small
\begin{tabular}{llll}
\toprule
 & (1) Baseline & (2) Entropy Premium & (3) Full \\
\midrule
H\_sentiment & $-0.0043^{**}$ &  &  \\
 & ($-2.39$) &  &  \\
Residual H &  & $+0.0008$ & $+0.0003$ \\
 &  & ($+0.17$) & ($+0.07$) \\
JS divergence & $-0.0001$ & $+0.0128$ & $+0.0114$ \\
 & ($-0.03$) & ($+1.06$) & ($+0.93$) \\
D\_post & $+0.0022$ & $-0.0049$ & $-0.0016$ \\
 & ($+0.96$) & ($-1.03$) & ($-0.43$) \\
H $\times$ D &  &  & $-0.0032$ \\
 &  &  & ($-1.15$) \\
\midrule
N months & 79 & 79 & 79 \\
\bottomrule
\end{tabular}
\end{table}
\noindent\textit{Notes:} Fama-MacBeth regressions with Newey-West lag = 6. All variables are standardized. Dependent variable: next-month stock return. ** $p<0.05$.

\subsection{Sub-Period Robustness}

Figure~\ref{fig:signals}, Panel B and Table~\ref{tab:subperiod} report sub-period FM regressions for residual H\_sentiment. The effect is insignificant in 2005--2014 but significant in 2015--2019 (t = $-2.37$) and 2020--2024 (t = $-2.77$). The increasing significance coincides with the rise of passive investing and ETF-driven herding, which may amplify the overconfidence channel: when market participants follow index flows rather than fundamental analysis, concentrated LLM sentiment is more likely to reflect (and reinforce) overconfidence.

\begin{table}[H]
\centering
\caption{Sub-Period Robustness: Residual H\_sentiment (US)}
\label{tab:subperiod}
\small
\begin{tabular}{lllll}
\toprule
 & 2005--2009 & 2010--2014 & 2015--2019 & 2020--2024 \\
\midrule
$\beta$ & $-0.013$ & $-0.021$ & $-0.079$ & $-0.110$ \\
$t$-stat & ($-0.37$) & ($-1.15$) & ($-2.37$) & ($-2.77$) \\
Sig. & & & ** & *** \\
N months & 20 & 20 & 20 & 20 \\
\bottomrule
\end{tabular}
\end{table}

\subsection{A-Share Cross-Validation}

Figure~\ref{fig:signals}, Panel C and Table~\ref{tab:ashare} report A-share FM regressions. H\_sentiment is significant (t = $-2.25$), JS divergence is significant (t = $+2.55$), and residual H\_sentiment is directionally consistent (t = $-1.80$). The A-share JS divergence result is notable: in a market with higher retail participation and stronger short-sale constraints, distributional distance between agents is more informative, consistent with the \citet{Miller1977} overpricing mechanism.

\begin{table}[H]
\centering
\caption{A-Share Cross-Validation (CSI 300, 2005--2024)}
\label{tab:ashare}
\small
\begin{tabular}{lllll}
\toprule
 & H\_sentiment & JS divergence & D\_post & Residual H \\
\midrule
$\beta$ & $-0.054$ & $+0.063$ & $+0.022$ & $-0.040$ \\
$t$-stat & ($-2.25$) & ($+2.55$) & ($+1.24$) & ($-1.80$) \\
Sig. & ** & ** & & * \\
N months & 80 & 80 & 80 & 80 \\
\bottomrule
\end{tabular}
\end{table}

\section{Robustness}\label{sec:robust}

\subsection{Decile Portfolio Sorts}

Figure~\ref{fig:decile} shows decile portfolios sorted on JS divergence. D1 (low JS) earns 0.98\%/mo while D10 (high JS) earns 2.44\%/mo, yielding a long-short return of 1.16\%/mo (t = 2.02, annualized Sharpe = 0.79). The JS decile sort is significant despite the FM regression being insignificant, suggesting a nonlinear relationship: extreme divergence is more informative than marginal divergence.

\begin{figure}[H]
\centering
\includegraphics[width=0.65\textwidth]{figures/fig_decile_js.png}
\caption{Decile Portfolio Returns Sorted by JS Divergence. Error bars represent $\pm$1.96 SE. Long-Short = D10 $-$ D1.}
\label{fig:decile}
\end{figure}

\subsection{Control Experiments}

We verify that the entropy premium is not an artifact of the LLM rating generation process. A quantitative model (without LLM) produces ratings with near-zero cross-sectional dispersion (mean std = 0.01), confirming that LLM agents generate genuine disagreement. The correlation between LLM and quantitative model ratings is 0.12, indicating that the two approaches capture different information.

\subsection{Multiple Testing Adjustment}

With 8 tests (4 signals $\times$ 2 markets), the Bonferroni-adjusted significance threshold is $0.05/8 = 0.00625$. Residual H\_sentiment (US) survives this adjustment ($p < 0.001$), while H\_sentiment (US, $p = 0.012$) and JS divergence (A-share, $p = 0.013$) do not. This underscores the importance of the orthogonality decomposition: the residual captures the independent component of entropy that is robust to multiple testing.

\section{Discussion}\label{sec:discussion}

\subsection{Why Entropy Dominates Disagreement}

The dominance of H\_sentiment over JS divergence and D\_post has a clear information-theoretic interpretation. Shannon entropy captures the \textit{full distribution} of agent sentiment, including higher-order moments that JS and D\_post miss. When agents agree on direction but disagree on intensity, D\_post is high but H\_sentiment is low (concentrated sentiment = overconfidence). When agents disagree on direction but cluster around similar intensity, D\_post is also high but H\_sentiment is high (dispersed sentiment = ambiguity). Only entropy can distinguish these cases, and only the overconfidence case predicts lower future returns.

Figure~\ref{fig:quasi} illustrates the quasi-experiment testing Proposition~\ref{prop:interaction}. The ``overconfidence'' portfolio (high H + high D) earns lower returns than the ``concordant'' portfolio (high H + low D), confirming that disagreement amplifies the overconfidence signal when sentiment is concentrated.

\begin{figure}[H]
\centering
\includegraphics[width=0.65\textwidth]{figures/fig_quasi_experiment.png}
\caption{Entropy $\times$ Disagreement Quasi-Experiment. Categories defined by standardized H\_sentiment and D\_post thresholds ($|z| > 0.5$). Testing Proposition~\ref{prop:interaction}: the H $\times$ D interaction captures the overconfidence premium.}
\label{fig:quasi}
\end{figure}

\subsection{LLM Agents as Information Environment Proxies}

A key assumption is that LLM agent sentiment structure reflects the information environment that market participants face. This assumption is plausible for three reasons. First, LLM agents process the same public information (earnings reports, news, analyst estimates) that market participants use. Second, the differentiation across agents (sentiment, technical, fundamental) mirrors the differentiation across market participants (momentum traders, value investors, sentiment-driven traders). Third, the entropy of LLM agent sentiment is correlated with observable measures of information concentration (analyst consensus, ETF ownership concentration), which we leave for future work.

\subsection{Implications for Multi-Agent AI System Design}

Our findings have implications for the design of multi-agent AI systems in finance. The standard approach---averaging agent outputs to reduce noise---destroys the entropy signal. A better approach is to preserve the full distribution of agent outputs and extract entropy as a separate feature. This is analogous to the insight in \citet{Chen2025} that inner confidence (token-level entropy) is a more informative signal than the point estimate.

\section{Conclusion}\label{sec:conclusion}

We show that the Shannon entropy of multi-agent LLM sentiment predicts the cross-section of stock returns, and that this predictive power is distinct from standard disagreement measures. The ``entropy premium''---the component of H\_sentiment orthogonal to JS divergence and D\_post---is the strongest and most robust signal (t = $-3.30$ in univariate FM, surviving Bonferroni adjustment). The effect is concentrated in the 2015--2024 period, consistent with the rise of passive investing amplifying overconfidence, and is directionally confirmed in A-share cross-validation.

Three theoretical propositions organize the empirical findings. Proposition~\ref{prop:entropy} establishes the entropy premium: concentrated sentiment signals overconfidence and predicts lower returns. Proposition~\ref{prop:distinct} establishes that entropy is distinct from disagreement: the orthogonal component is more informative than the raw measure. Proposition~\ref{prop:interaction} establishes the entropy--disagreement interaction: overconfidence is strongest when sentiment is concentrated AND disagreement is high.

These findings have two policy implications. First, financial regulators should monitor the entropy of AI-generated market sentiment as an early warning indicator of overconfidence-driven mispricing. Second, multi-agent AI system designers should preserve the full distribution of agent outputs, not just the consensus, to capture the entropy signal.

\bibliographystyle{plainnat}
\bibliography{references}

\end{document}
